\begin{figure}[H]
    \centering
    \tdplotsetmaincoords{70}{135}
    \begin{tikzpicture}[
        tdplot_main_coords,
        scale=1.45,
        line cap=round,
        line join=round
    ]
        % Pontos especiais do grafico.
        \coordinate (origem) at (0,0,0);
        \coordinate (pontoFaltante) at (0,0,1);

        % Projecao do plano z=1 no dominio R^2.
        \draw[gray!60, dashed] (-2.25,-1.60,0) -- (2.25,-1.60,0) --
            (2.25,1.60,0) -- (-2.25,1.60,0) -- cycle;

        % Eixos coordenados, com o eixo vertical representando os valores de f.
        \draw[->, very thick] (-2.75,0,0) -- (2.95,0,0) node[below] {$x$};
        \draw[->, very thick] (0,-2.25,0) -- (0,2.55,0) node[right] {$y$};
        \draw[very thick] (0,0,-0.30) -- (0,0,1);



        % Parte principal do grafico: plano z=1 para (x,y) diferente de (0,0).
        \path[fill=Cerulean!14, fill opacity=0.55, draw=Cerulean!70!black, thick]
            (-2.25,-1.60,1) -- (2.25,-1.60,1) --
            (2.25,1.60,1) -- (-2.25,1.60,1) -- cycle;



        % Segmento vertical sobre a origem, indicando a descontinuidade.
        \draw[gray!65, dashed] (origem) -- (pontoFaltante);

        % Trecho superior do eixo z, desenhado apos o plano para ficar opaco.
        \draw[->, very thick] (0,0,1) -- (0,0,1.95) node[left] {$f(x,y)=z$};

        % Rotulo do ponto que pertence ao grafico.
        \node[black, right=7pt] at (origem) {$(0,0,f(0,0))$};

        % Ponto aberto: o valor 1 seria o limite, mas nao e f(0,0).
        \draw[Cerulean!70!black, very thick, fill=white] (pontoFaltante) circle (2.8pt);

        % Rotulo do ponto removido do plano z=1.
        \node[black, right=7pt] at (pontoFaltante) {$(0,0,1)$};
        % Valor real da funcao na origem: f(0,0)=0.
        \filldraw[Cerulean!70!black] (origem) circle (2.5pt);
    \end{tikzpicture}
    \caption{Gráfico de $f$, com eixo vertical $z=f(x,y)$.}
    \label{fig:limiteNaoImporta}
\end{figure}
